% =============================================================================
% HISTORIAL DE REVISIONES Y ROLES
% =============================================================================
\chapter*{Historial de Revisiones y Roles}
\addcontentsline{toc}{chapter}{Historial de Revisiones y Roles}
\markboth{Historial de Revisiones y Roles}{}

El presente documento constituye un producto de información de usuario del ciclo de
vida del software conforme al estándar \textbf{ISO/IEC/IEEE 26514:2022}, que define el
diseño y desarrollo de la documentación de usuario para sistemas y software. Su
elaboración, revisión y aprobación quedan trazadas en los registros siguientes.

\section*{Control de Versiones del Documento}

\vspace{0.2cm}
\noindent
\renewcommand{\arraystretch}{1.4}
\fontsize{9}{10.8}\selectfont\sffamily
\begin{tabularx}{\textwidth}{|c|c|Y|c|}
\hline
\rowcolor{headerTabla}
\textbf{Versión} & \textbf{Fecha} & \textbf{Descripción del cambio} & \textbf{Estado} \\
\hline
0.1.0 & 2026-05-02 & Estructura inicial del manual y definición del índice maestro alineado a ISO/IEC/IEEE 26514:2022. & Borrador \\
\hline
\rowcolor{grisTecnico}
0.6.0 & 2026-05-28 & Documentación operativa de los perfiles Profesional y Reclutador con apoyo visual paso a paso. & En revisión \\
\hline
1.0.0 & 2026-06-18 & Documentación completa de los tres perfiles (Profesional, Reclutador y Administrador), diseño adaptable, resolución de problemas y glosario. & Aprobado \\
\hline
\end{tabularx}

\vspace{0.6cm}

\clearpage
\section*{Roles del Equipo de Documentación}
\addcontentsline{toc}{section}{Roles del Equipo de Documentación}
La elaboración de este manual fue responsabilidad del equipo de \textbf{Byte Busters
S.R.L.}, con la asignación de roles y validaciones que se detalla a continuación.

\vspace{0.2cm}
\noindent
\renewcommand{\arraystretch}{1.4}
\fontsize{9}{10.8}\selectfont\sffamily
\begin{tabularx}{\textwidth}{|>{\bfseries\color{colorPrimario}}l|Y|c|}
\hline
\rowcolor{headerTabla}
\textbf{Rol} & \textbf{Responsabilidad en el documento} & \textbf{Validación} \\
\hline
Líder de Documentación & Definición de la estructura, estilo editorial y coherencia con el estándar 26514. & Editorial \\
\hline
\rowcolor{grisTecnico}
Analista de Producto & Redacción de los procedimientos paso a paso y captura de las pantallas del sistema. & Funcional \\
\hline
Especialista UX & Revisión de la nomenclatura de controles, accesibilidad y claridad de las instrucciones. & Usabilidad \\
\hline
\rowcolor{grisTecnico}
QA Lead & Verificación de la correspondencia entre los procedimientos descritos y el comportamiento real del sistema. & Calidad \\
\hline
\end{tabularx}

\vspace{0.4cm}
\begin{NotaTecnica}
Este manual está dirigido a las personas usuarias finales de la plataforma EthosHub.
Su identificador de trazabilidad es \comando{MU-PERF-ETHOSHUB-1.0.0}. Las capturas de
pantalla corresponden a la versión vigente del sistema y pueden variar ligeramente
según las actualizaciones de la interfaz.
\end{NotaTecnica}
